% Options for packages loaded elsewhere
\PassOptionsToPackage{unicode}{hyperref}
\PassOptionsToPackage{hyphens}{url}
\documentclass[
]{article}
\usepackage{xcolor}
\usepackage[margin=1in]{geometry}
\usepackage{amsmath,amssymb}
\setcounter{secnumdepth}{5}
\usepackage{iftex}
\ifPDFTeX
  \usepackage[T1]{fontenc}
  \usepackage[utf8]{inputenc}
  \usepackage{textcomp} % provide euro and other symbols
\else % if luatex or xetex
  \usepackage{unicode-math} % this also loads fontspec
  \defaultfontfeatures{Scale=MatchLowercase}
  \defaultfontfeatures[\rmfamily]{Ligatures=TeX,Scale=1}
\fi
\usepackage{lmodern}
\ifPDFTeX\else
  % xetex/luatex font selection
\fi
% Use upquote if available, for straight quotes in verbatim environments
\IfFileExists{upquote.sty}{\usepackage{upquote}}{}
\IfFileExists{microtype.sty}{% use microtype if available
  \usepackage[]{microtype}
  \UseMicrotypeSet[protrusion]{basicmath} % disable protrusion for tt fonts
}{}
\makeatletter
\@ifundefined{KOMAClassName}{% if non-KOMA class
  \IfFileExists{parskip.sty}{%
    \usepackage{parskip}
  }{% else
    \setlength{\parindent}{0pt}
    \setlength{\parskip}{6pt plus 2pt minus 1pt}}
}{% if KOMA class
  \KOMAoptions{parskip=half}}
\makeatother
\usepackage{longtable,booktabs,array}
\usepackage{calc} % for calculating minipage widths
% Correct order of tables after \paragraph or \subparagraph
\usepackage{etoolbox}
\makeatletter
\patchcmd\longtable{\par}{\if@noskipsec\mbox{}\fi\par}{}{}
\makeatother
% Allow footnotes in longtable head/foot
\IfFileExists{footnotehyper.sty}{\usepackage{footnotehyper}}{\usepackage{footnote}}
\makesavenoteenv{longtable}
\usepackage{graphicx}
\makeatletter
\newsavebox\pandoc@box
\newcommand*\pandocbounded[1]{% scales image to fit in text height/width
  \sbox\pandoc@box{#1}%
  \Gscale@div\@tempa{\textheight}{\dimexpr\ht\pandoc@box+\dp\pandoc@box\relax}%
  \Gscale@div\@tempb{\linewidth}{\wd\pandoc@box}%
  \ifdim\@tempb\p@<\@tempa\p@\let\@tempa\@tempb\fi% select the smaller of both
  \ifdim\@tempa\p@<\p@\scalebox{\@tempa}{\usebox\pandoc@box}%
  \else\usebox{\pandoc@box}%
  \fi%
}
% Set default figure placement to htbp
\def\fps@figure{htbp}
\makeatother
\setlength{\emergencystretch}{3em} % prevent overfull lines
\providecommand{\tightlist}{%
  \setlength{\itemsep}{0pt}\setlength{\parskip}{0pt}}
\usepackage[]{natbib}
\bibliographystyle{plainnat}
\usepackage{booktabs}
\usepackage{longtable}
\usepackage{array}
\usepackage{multirow}
\usepackage{wrapfig}
\usepackage{float}
\usepackage{colortbl}
\usepackage{pdflscape}
\usepackage{tabu}
\usepackage{threeparttable}
\usepackage{threeparttablex}
\usepackage[normalem]{ulem}
\usepackage{makecell}
\usepackage{xcolor}
\usepackage{bookmark}
\IfFileExists{xurl.sty}{\usepackage{xurl}}{} % add URL line breaks if available
\urlstyle{same}
\hypersetup{
  pdftitle={Título del Trabajo Final},
  pdfauthor={Isabella Palacios; Mayerlin Otamendi},
  hidelinks,
  pdfcreator={LaTeX via pandoc}}

\title{Título del Trabajo Final}
\author{Isabella Palacios \and Mayerlin Otamendi}
\date{Abril 2026}

\begin{document}
\maketitle
\begin{abstract}
En el presente informe se realiza un análisis comparativo sobre el
impacto de la pandemia en el consumo de alcohol per cápita en Europa. Se
examina la tendencia del consumo de alcohol en la región europea antes
(2016-2019), durante (2020) y después (2022) del COVID-19 para
identificar las variaciones causadas por el mismo, aplicando técnicas de
estadística descriptiva y visualización de datos con R.
\end{abstract}

{
\setcounter{tocdepth}{2}
\tableofcontents
}
\section{Introducción}

La región europea ha mantenido históricamente los niveles más altos de
consumo de alcohol per cápita en el mundo. Entre 2016 y 2019, la
tendencia se caracterizó por una relativa estabilidad en la tendencia de
consumo social, fuertemente arraigados en la cultura y la economía del
continente. Sin embargo, la aparición del COVID-19 a finales de 2019 y
su declaración como pandemia por la OMS en marzo de 2020 transformaron
radicalmente este panorama.

Durante el año 2020, las restricciones de movilidad y los confinamientos
obligatorios en Europa provocaron el cierre masivo de la hostelería
(bares, restaurantes y clubes). No obstante, la consideración del
alcohol como un ``bien esencial'' en muchos países europeos permitió que
los puntos de venta minorista y el comercio electrónico mantuvieran la
disponibilidad del producto.

Para el año 2022, con la flexibilización de las medidas sanitarias y la
reapertura de los espacios públicos, el comportamiento del consumidor
europeo entró en una fase de reajuste. Este periodo de ``postpandemia''
es crítico para analizar si las variaciones en la frecuencia de consumo
observados durante el aislamiento se consolidaron como hábitos crónicos
o si regresaron a los niveles base de 2016-2019.

En la presente investigación se busca analizar cómo se comportan los
niveles del consumo de alcohol en el continente europeo antes
(2016-2019), durante (2020) y después (2022) de la pandemia del
COVID-19, analizando la varianza del consumo de alcohol per cápita entre
los países y comparar sus promedios postpandemia con el promedio del
consumo antes de la misma.

\section{Planteamiento del Problema}

Describir planteamiento del problema\\

La región europea se ha distinguido históricamente por registrar los
niveles mas elevados de consumo de alcohol per cápita a escala global.
Hasta 2019, este consumo presentaba una tendencia a la estabilidad,
integrada profundamente en las dinámicas sociales y económicas del
continente, sin embargo, la irrupción de la pandemia de COVID-19 en 2020
generó una disrupción sin precedentes en estos patrones.

El problema radica en la incertidumbre sobre la naturaleza de los
cambios observados durante la crisis sanitaria. Durante el
confinamiento, el cierre de la hostelería desplazó el consumo hacia el
hogar, facilitado por el comercio electrónico y la clasificación del
alcohol como ``bien esencial''. Si bien algunos sectores de la población
redujeron su ingesta, otros la aumentaron como mecanismo de
afrontamiento ante el estrés y el aislamiento.

La falta de un análisis comparativo integral que contraste el periodo
base (2016-2019) con el año de crisis (2020) y la recuperación (2022)
impide identificar qué países presentan comportamientos extremos y
cuáles han logrado retornar a la normalidad previa. Esta brecha de
información dificulta la creación de políticas de salud pública
adaptadas a la nueva realidad europea. Debido a esto, se plantea la
siguiente cuestión,

\textbf{Pregunta de investigacón: ¿Estamos ante una anomalía estadística provocada por las restricciones de movilidad o ante una transformación a largo plazo en la relación del ciudadano europeo con el consumo de alcohol?}

\section{Justificación}

La justificación de este trabajo de investigación reside en la necesidad
de determinar el impacto de la pandemia en cuanto al consumo de alcohol
en el continente europeo, al no entender claramente como afectó esta
crisis en las dinámicas de consumo no se puede crear una estrategia de
salud publica coherente con la realidad actual para evitar excesos, y
por otro lado el mercado del alcohol en Europa necesita estar alineado a
los nuevos patrones de consumo para garantizar sus ventas.

Desde el punto de vista académico la realización de este informe
significa un aporte para otros investigadore debido a que cuenta con una
metodología simple y replicable con el uso de herramientas de fácil
manejo y acceso, siendo la principal R Markdown.

\section{Objetivos}

\subsection{Objetivo General}
\begin{itemize}
  \item{Evaluar de manera comparativa los patrones de consumo de alcohol en la población europea entre los años 2016 y 2022 analizando el impacto de la pandemia de COVID-19 para identificar cambios en la ingesta antes, durante y después de la crisis sanitaria.}
\end{itemize}

\subsection{Objetivos Específicos}
\begin{itemize}
  \item{Realizar un análisis de estadística descriptiva sobre el consumo de alcohol per cápita para comparar el comportamiento de la región europea en los periodos 2016-2019, 2020 y 2022.}
  \item{Comparar el consumo de alcohol de cada país en 2020 y 2022 frente al promedio del periodo 2016-2019, para determinar si las variaciones derivadas de la pandemia fueron transitorias o persistentes.}
  \item{Evaluar la homogeneidad de los datos a través del Coeficiente de Variación, con el fin de identificar qué tan representativo es el promedio regional en cada temporada de la investigación.}
  \item{Identificar y analizar los datos atípicos del consumo de alcohol en 2020 y 2022 para señalar que países presentaron comportamientos extremos durante y después de la pandemia.}
\end{itemize}

\section{Cobertura}

\subsection{Horizontal}

La presente investigación delimita su cobertura horizontal a los países
miembros del continente europeo analizando los reportes de datos
oficiales sobre el consumo de alcohol per cápita (medido en litros de
alcohol puro al año). No solo se analiza el volumen total, sino también
la variabilidad estadística (homogeneidad) y la detección de
comportamientos extremos (datos atípicos).

\subsection{Vertical}

En cuanto a su cobertura vertical, el análisis profundiza en la
evolución temporal desde 2016 hasta 2022, dividiendo este periodo en
tres fases para la comparación: \textbf{(2016-2019)} comportamiento
pre-pandemia, \textbf{(2020)} impacto inmediato de las restricciones
sanitarias y \textbf{(2022)} determina si los cambios fueron
transitorios o persistentes.

\section{Periodo de Referencia}

\section{Variables}

Las principales variables en estudio son ``Litros de alcohol per
cápita'', que representa la cantidad promedio de alcohol puro consumido
por persona al año, y ``Año'' que abarca el periodo prepandemia
(2016-2019), durante la pandemia (2020) y pospandemia (2022). Ambas son
dependientes entre sí de forma que el análisis conjunto es crucial para
responder a los objetivos de la investigación. No menos importante
tenemos a la variable ``País'' que abarca a los países de la región
europea de los cuales se analizan los niveles de consumo de alcohol
dependiendo del período que se esté tratando, para luego hacer una
comparación que dé respuesta al objetivo de la investigación.

\section{Marco Teórico}

\section{Bases Teóricas}

El presente estudio se sustenta en la estadística descriptiva, la cual
permite organizar, resumir y presentar los datos de manera informativa
para identificar patrones en el consumo de alcohol en Europa. Para dar
cumplimiento a los objetivos específicos y respuesta a la pregunta de
investigación, se implementa un análisis de variación porcentual y
detección de datos atípicos haciendo uso de las siguientes medidas
estadísticas:

\begin{itemize}
\item
  \textbf{Esperanza Matemática (Media Muestral):}
  \[\bar{x} = \frac{1}{n} \sum_{i=1}^{n} x_i\]
\item
  \textbf{Varianza (Volatilidad):}
  \[S^2 = \frac{1}{n-1} \sum_{i=1}^{n} (x_i - \bar{x})^2\]
\item
  \textbf{Desviación Estándar:}
  \[s = \sqrt{\frac{1}{n-1} \sum_{i=1}^{n} (x_i - \bar{x})^2}\]
\end{itemize}

\subsection{El Consumo de Alcohol en Europa como Objeto de Estudio}

El consumo de alcohol no es un fenómeno estático sino que responde a
dinámicas de adaptación ante el estrés y el aislamiento social.
Históricamente, la región europea ha registrado los niveles más altos de
consumo de alcohol per cápita a nivel mundial. Este fenómeno no es solo
una estadística de salud pública, sino que está profundamente arraigado
en la cultura, las interacciones sociales y la economía del
continente.\\

\subsection{Impacto dela Pandemia de COVID-19}

Entre 2016 y 2019, la tendencia se caracterizó por una relativa
estabilidad en los patrones de consumo social. La aparición del COVID-19
a finales de 2019 у su posterior declaración como pandemia por la OMS en
marzo de 2020 transformaron radicalmente este panorama. Las
restricciones de movilidad y los confinamientos obligatorios provocaron
el cierre masivo del sector de la hostelería, incluyendo bares,
restaurantes y clubes. No obstante, en muchos países europeos, el
alcohol fue considerado un ``bien esencial'', lo que permitió que los
puntos de venta minorista y el comercio electrónico mantuvieran la
disponibilidad del producto.\\

\section{Marco Metódico}

En esta sección se presentan las metodologías y/o test necesarios
relacionados con la práctica divulgada en el marco teórico, así como las
fuentes de datos.

\section{Bases metódicas utilizadas en el trabajo de investigación}

Desarrollo\\

\subsection{Item 1}

Descripción\\

\subsection{Item 2}

Descripción\\

\subsubsection{Item 2.2}

Descripción\\

\subsection{Item n}

Descripción\\

\section{Análisis de Resultados}

\{\small Este capítulo presenta los resultados obtenidos al aplicar las
metodologías descritas en el capítulo anterior, así como del análisis
estadístico descriptivo de las variables estudiadas y las ventajas y
desventajas al aplicar dichos tests.\}

\section{Presentación de resultados 1}

Descripción\\

\begin{longtable}[]{@{}
  >{\raggedright\arraybackslash}p{(\linewidth - 8\tabcolsep) * \real{0.3333}}
  >{\raggedleft\arraybackslash}p{(\linewidth - 8\tabcolsep) * \real{0.2133}}
  >{\raggedleft\arraybackslash}p{(\linewidth - 8\tabcolsep) * \real{0.1333}}
  >{\raggedleft\arraybackslash}p{(\linewidth - 8\tabcolsep) * \real{0.1200}}
  >{\raggedleft\arraybackslash}p{(\linewidth - 8\tabcolsep) * \real{0.2000}}@{}}
\caption{Análisis Descriptivo del Consumo de Alcohol Per Cápita en
Europa}\tabularnewline
\toprule\noalign{}
\begin{minipage}[b]{\linewidth}\raggedright
Periodo
\end{minipage} & \begin{minipage}[b]{\linewidth}\raggedleft
N° de Registros
\end{minipage} & \begin{minipage}[b]{\linewidth}\raggedleft
Media (L)
\end{minipage} & \begin{minipage}[b]{\linewidth}\raggedleft
Varianza
\end{minipage} & \begin{minipage}[b]{\linewidth}\raggedleft
Desv. Estándar
\end{minipage} \\
\midrule\noalign{}
\endfirsthead
\toprule\noalign{}
\begin{minipage}[b]{\linewidth}\raggedright
Periodo
\end{minipage} & \begin{minipage}[b]{\linewidth}\raggedleft
N° de Registros
\end{minipage} & \begin{minipage}[b]{\linewidth}\raggedleft
Media (L)
\end{minipage} & \begin{minipage}[b]{\linewidth}\raggedleft
Varianza
\end{minipage} & \begin{minipage}[b]{\linewidth}\raggedleft
Desv. Estándar
\end{minipage} \\
\midrule\noalign{}
\endhead
\bottomrule\noalign{}
\endlastfoot
2016-2019 (Pre-Pandemia) & 164 & 10.12 & 5.73 & 2.39 \\
2020 (Pandemia) & 41 & 10.04 & 6.24 & 2.50 \\
2022 (Post-Pandemia) & 41 & 10.02 & 6.58 & 2.57 \\
\end{longtable}

\section{Presentación de resultados 2}

Descripción\\

\begin{center}\includegraphics[width=0.9\linewidth,]{Informe_Trabajo_Final_files/figure-latex/objetivo2-barras-1} \end{center}

\section{Presentación de resultados n}

Descripción\\

\begin{figure}[H]

{\centering \includegraphics[width=0.9\linewidth,]{Informe_Trabajo_Final_files/figure-latex/objetivo3-homogeneidad-1} 

}

\caption{Evolución de la homogeneidad de los datos (Coeficiente de Variación) por año.}\label{fig:objetivo3-homogeneidad}
\end{figure}

\section{Presentación de resultados n}

Descripción\\

\begin{center}\includegraphics[width=0.9\linewidth,]{Informe_Trabajo_Final_files/figure-latex/top-5-caidas-1} \end{center}

\section{Presentación de resultados n}

Descripción\\

\begin{center}\includegraphics[width=0.9\linewidth,]{Informe_Trabajo_Final_files/figure-latex/boxplot-identificado-1} \end{center}

\section{Conclusiones y Recomendaciones}

\textbf{1.-} Conclusion 1 del trabajo de investigación.\\

\textbf{Recomendación} Recomendación 1 del trabajo de investigación.\\

\textbf{2.-} Conclusion 2 del trabajo de investigación.\\

\textbf{Recomendación} Recomendación 2 del trabajo de investigación.\\

\textbf{n.-} Conclusion n del trabajo de investigación.\\

\textbf{Recomendación} Recomendación n del trabajo de investigación.\\

\section{Anexo A}

\section{Anexo B}

\section{Anexo C}

\bibliography{bibliografia.bib}

\end{document}
